\chapter{Design Philosophy}
\label{chap:design-philosophy}

CharlotteOS combines ideas from three projects around one thesis.

\section{The combined thesis}

\begin{quote}
\textbf{Hardware is asynchronous; the operating system and the userspace runtime
should be too.}
\end{quote}

Three systems approach the same design space from different directions. Rather than
collapsing their abstractions into one universal mechanism, the architecture
composes them at the right boundaries.

\subsection{CharlotteOS --- the protection, scheduling, and interrupt substrate}

CharlotteOS contributes the mechanisms that \emph{must} be privileged:

\begin{itemize}
    \item Address-space isolation through page tables and the MMU.
    \item Kernel scheduling of threads across logical processors.
    \item LP-local state and interrupt routing.
    \item Page-table manipulation and physical-memory management.
    \item MMIO and DMA authority delegation.
    \item Per-domain capability tables.
    \item Event observation and thread wakeup.
    \item Syscall dispatch and asynchronous completion delivery.
\end{itemize}

Its design rule is: \textbf{the kernel provides the minimum set of primitives from
which everything else can be composed.} If a facility can be built in userspace
without compromising isolation, it belongs in userspace.

\subsection{Sitas --- the shard-per-core execution discipline}

Sitas (\url{https://github.com/FrodeRanders/sitas}) contributes an execution discipline --- a set of rules about how state is
owned and mutated:

\begin{quote}
Only the owning shard mutates shard-local state. Cross-shard interaction occurs
through bounded typed messages carrying owned values.
\end{quote}

Sitas provides:

\begin{itemize}
    \item One executor per shard --- cooperative task scheduling within
        a shard; preemption between shards.
    \item Futures as userspace state machines --- every operation that may
        wait is a \texttt{Future}, driven by the shard's executor.
    \item Bounded backpressure --- every queue has fixed capacity; senders
        must handle \texttt{WouldBlock}.
    \item Shard-local ownership --- mutable state belongs to exactly one
        shard; no locks, no atomics for shard-internal data.
    \item Explicit placement --- where a task runs is a deliberate choice,
        not an accident of scheduling.
    \item Typed command/reply APIs --- what a shard receives is a specific
        Rust type, not a generic byte buffer.
    \item Structured shutdown and cancellation --- tasks can be cancelled,
        and cancellation has defined ownership semantics.
    \item Observability through owned snapshots --- diagnostics capture
        state by value, never by borrowed reference into runtime internals.
\end{itemize}

The sitas discipline is enforced by Rust's type system within a single protection
domain. CharlotteOS extends it across domains through the kernel.

\subsection{Xous --- the isolated userspace service model}

Xous, a microkernel OS for Precursor and Betrusted hardware, contributes the
protection-domain service model:

\begin{itemize}
    \item Services are userspace servers --- drivers, filesystems, and
        protocol stacks run in their own protection domains, not in the kernel.
    \item Clients connect to servers and receive connection identifiers ---
        a connection \emph{is} authority; possessing one means you may communicate
        with that server.
    \item IPC is a kernel primitive --- communication is not a convention
        layered on shared memory; it is mediated and validated by the kernel.
    \item Scalar messages are distinct from memory-bearing messages ---
        small control-plane operations (opcodes, handles, status codes) travel
        in fixed-size messages; bulk data travels as memory-object attachments.
    \item Memory IPC distinguishes ownership transfer from borrowing ---
        move transfers ownership permanently. Borrow grants temporary access,
        revoked on reply or cancellation.
    \item A userspace name service maps human-readable names to opaque
        server identities --- the kernel knows nothing about strings.
\end{itemize}

CharlotteOS adopts Xous's userspace-server model and combines it with the sitas
execution discipline and an async-native kernel. The detailed co-design is
maintained in \repofile|docs/architecture/sitas-xous.md|.

\section{How they fit together}

The three systems do not collapse into one. They compose at well-defined
boundaries:

\begin{center}
\begin{tabular}{|l|l|}
\hline
\textbf{Layer} & \textbf{Provided by} \\
\hline
Shard-local async execution & sitas \\
Cross-domain protection and IPC & CharlotteOS kernel (Xous-inspired) \\
Completion delivery for kernel/device ops & CharlotteOS kernel \\
Service discovery and policy & Userspace (name service, supervisor) \\
\hline
\end{tabular}
\end{center}

The layers are:

\begin{center}
\begin{tikzpicture}[
    layer/.style={draw, rounded corners, align=left, text width=11cm,
        inner sep=7pt, font=\small},
]
    \node[layer] (client) {%
        \textbf{Client process}\par\vspace{2pt}
        \begin{itemize}
            \setlength{\itemsep}{0pt}
            \item typed protocol library
            \item endpoint connection capabilities
            \item futures backed by IPC replies or completion records
            \item owned buffers whose transfer mode is explicit
        \end{itemize}};

    \node[layer, below=0.4cm of client] (service) {%
        \textbf{Userspace service process}\par\vspace{2pt}
        \begin{itemize}
            \setlength{\itemsep}{0pt}
            \item one or more externally visible IPC endpoints
            \item protocol-specific request/reply handling
            \item one sitas executor per assigned shard
            \item shard-local state
            \item typed in-process commands
            \item driver or service implementation
        \end{itemize}};

    \node[layer, below=0.4cm of service] (kernel) {%
        \textbf{CharlotteOS kernel}\par\vspace{2pt}
        \begin{itemize}
            \setlength{\itemsep}{0pt}
            \item protection domains
            \item threads and LP scheduling
            \item capabilities and endpoint connections
            \item interrupt routing
            \item memory objects and page mappings
            \item operation submission and completion queues
        \end{itemize}};
\end{tikzpicture}
\end{center}

\architecturefigure{system-layering}
    {CharlotteOS layering from hardware and kernel mechanism through isolated
    services and generated application logic to managed external systems.}
    {fig:system-layering}

\section{The three kernel primitives}

Three orthogonal kernel abstractions form the foundation:

\begin{center}
\begin{tikzpicture}[
    box/.style={draw, rounded corners, align=center, minimum width=2.8cm, inner sep=3pt},
    note/.style={align=center, font=\footnotesize},
]
    \node[box] (root) {Kernel Objects};
    \node[box, below left=1.5cm and 2.1cm of root] (caps) {Capabilities};
    \node[box, below=1.5cm of root] (ends) {Endpoints};
    \node[box, below right=1.5cm and 2.1cm of root] (mems) {Memory Objects};

    \node[note, below=0.3cm of caps] {Authority\\Security};
    \node[note, below=0.3cm of ends] {Communication\\Rendezvous};
    \node[note, below=0.3cm of mems] {Ownership\\Data Movement};

    \coordinate (mid) at ([yshift=-0.75cm]root.south);
    \draw (root.south) -- (mid);
    \draw (caps.north |- mid) -- (mems.north |- mid);
    \draw (caps.north) -- (caps.north |- mid);
    \draw (ends.north) -- (ends.north |- mid);
    \draw (mems.north) -- (mems.north |- mid);
\end{tikzpicture}
\end{center}

Every other kernel facility --- completion queues, reply tokens, device grants,
DMA mappings --- is a \textbf{composition} of these three primitives.

\subsection{Capabilities: ``May I?''}

A capability answers one question: \textbf{``May I perform this operation?''}
It represents authority, not work. Examples include endpoint capabilities,
connection capabilities, memory-object capabilities, interrupt capabilities,
DMA-domain capabilities, and timer capabilities.

\subsection{Endpoints: ``With whom?''}

An endpoint answers: \textbf{``Who am I communicating with?''} It provides
bounded rendezvous between protection domains. Endpoints know nothing about packet
buffers, futures, DMA, or ownership. They carry requests and replies.

\subsection{Memory Objects: ``Where is the data, and who owns it?''}

A memory object answers: \textbf{``Where is the data, and who currently owns
it?''} It represents page-backed storage whose ownership can move independently of
communication. Its responsibilities include page ownership, mappings, transfer
modes, DMA state, and lifetime.

\section{An explicit design rule}

\begin{quote}
\textbf{When adding a new subsystem, first ask whether it can be expressed using
capabilities, endpoints, and memory objects. New kernel primitives should be
introduced only when they cannot be naturally composed from the existing ones.}
\end{quote}

This rule limits the kernel interface and keeps the ABI and trust model reviewable.

\section{What ``asynchronous throughout'' means}

``Asynchronous throughout'' means:

\begin{enumerate}
    \item Operations that may wait for external progress have submission/completion
        semantics.
    \item No subsystem boundary forces the caller to dedicate a thread to waiting
        for one operation.
    \item A userspace shard can wait on one aggregated kernel source (the CQ).
    \item Completions are retained until observed.
    \item Buffers and capabilities have explicit ownership during an operation.
    \item Cancellation and teardown have defined state transitions.
    \item Bounded queues preserve backpressure across layers.
    \item Service implementations can suspend one task while continuing other work
        on the shard.
\end{enumerate}

This is not the same as saying:

\begin{itemize}
    \item Every syscall is asynchronous.
    \item Every call returns a newly allocated completion capability.
    \item No thread ever blocks.
    \item Every wake sends an IPI.
    \item All waitable objects have identical semantics.
    \item Userspace code is interrupted by arbitrary asynchronous callbacks.
\end{itemize}

\subsection{Blocking is still useful}

When no sitas task is runnable, the shard's execution context should block in one
kernel wait operation:

\begin{center}
\begin{tikzpicture}[
    cond/.style={draw, dashed, rounded corners, align=center, minimum width=2.8cm, inner sep=4pt},
    act/.style={draw, rounded corners, align=center, inner sep=4pt},
]
    \node[cond] (ready) {ready tasks\\exist};
    \node[act, right=0.9cm of ready] (poll) {poll tasks};

    \node[cond, below=0.6cm of ready] (noready) {no ready\\tasks};
    \node[act, right=0.9cm of noready] (wait) {wait on CQ /\\endpoint readiness /\\deadline};

    \node[cond, below=0.6cm of noready] (event) {event\\arrives};
    \node[act, right=0.9cm of event] (runnable) {execution context\\becomes runnable};
    \node[act, right=0.7cm of runnable] (drain) {drain events};
    \node[act, right=0.7cm of drain] (wake) {wake relevant\\futures};

    \draw[->, dashed] (ready) -- (poll);
    \draw[->, dashed] (noready) -- (wait);
    \draw[->, dashed] (event) -- (runnable);
    \draw[->] (runnable) -- (drain);
    \draw[->] (drain) -- (wake);
\end{tikzpicture}
\end{center}

\begin{quote}
\textbf{The shard parks directly on native completion and message state; no
helper-thread pool is required merely to convert blocking kernel interfaces
into futures.}
\end{quote}

\subsection{Immediate operations remain immediate}

Operations that cannot meaningfully block --- capability duplication, metadata
queries, nonblocking CQ drain, closing a quiescent object --- complete
synchronously. The ABI does not force asynchronous ceremony onto inherently
synchronous work.

\subsection{No asynchronous userspace upcalls}

Completions do not inject callbacks into userspace at arbitrary instructions. The
path is:

\begin{center}
\begin{tikzpicture}[
    cond/.style={draw, dashed, rounded corners, align=center, inner sep=5pt},
    stage/.style={draw, rounded corners, align=center, inner sep=5pt},
]
    \node[cond] (s1) {hardware interrupt};
    \node[stage, below=0.4cm of s1] (s2) {kernel records result\\or notification};
    \node[stage, below=0.4cm of s2] (s3) {blocked execution context\\becomes runnable};
    \node[stage, below=0.4cm of s3] (s4) {normal return to userspace};
    \node[stage, below=0.4cm of s4] (s5) {sitas executor drains\\CQ/endpoints};
    \node[stage, below=0.4cm of s5] (s6) {executor wakes\\relevant Rust tasks};
    \draw[->] (s1) -- (s2);
    \draw[->] (s2) -- (s3);
    \draw[->] (s3) -- (s4);
    \draw[->] (s4) -- (s5);
    \draw[->] (s5) -- (s6);
\end{tikzpicture}
\end{center}

The kernel may preempt threads, but it does not inject arbitrary task callbacks
into userspace. Task scheduling therefore remains cooperative and reviewable.

\section{Why this matters for critical software}

Every design decision in CharlotteOS traces back to one or more of the six
principles from \autoref{chap:problem}:

\begin{itemize}
    \item Capabilities replace ambient authority --- a compromised service
        cannot access resources it was never granted. The set of things a process
        can do is enumerable and auditable by construction.
    \item Endpoints make communication authenticated --- a client that holds
        a connection capability has proven (at grant time) that it is authorized
        to communicate. The kernel enforces this on every message.
    \item Memory objects constrain stale cross-domain access --- ownership moves
        through the kernel and is enforced by page tables. Correctness still
        depends on MMU, TLB, and DMA handling.
    \item Composition limits the privileged code footprint --- most system functionality
        lives in userspace, where faults can be contained to a service domain.
    \item The async model reduces adaptation --- the kernel and runtime share
        submission, completion, and wake semantics.
    \item Bounded ingress makes pressure visible --- endpoint and hardware queues
        apply explicit backpressure. \autoref{chap:scheduling} records the remaining unbounded
        internal backlog.
\end{itemize}

\endinput
